\documentclass[11pt]{article}

\usepackage[margin=1in]{geometry}
\usepackage[T1]{fontenc}
\usepackage{lmodern}
\usepackage{booktabs}
\usepackage{enumitem}
\usepackage{xcolor}
\usepackage{listings}
\usepackage[hidelinks]{hyperref}

\definecolor{codebg}{gray}{0.96}
\lstset{
  language=Python,
  basicstyle=\ttfamily\small,
  backgroundcolor=\color{codebg},
  frame=single,
  rulecolor=\color{gray!40},
  columns=fullflexible,
  keepspaces=true,
  showstringspaces=false,
  keywordstyle=\bfseries,
  commentstyle=\itshape\color{gray!70!black},
  xleftmargin=4pt,
  xrightmargin=4pt,
}
\setlist{itemsep=2pt, topsep=4pt}

\title{Project 4: Jellyfish\\[4pt] \large Simulator Guide for Students}
\author{COMS 4444}
\date{}

\begin{document}
\maketitle
\tableofcontents

% =====================================================================
\section{How to code your jellyfish}

\subsection{Step by step}

\begin{enumerate}[leftmargin=*]
  \item \textbf{Open your group's file.} Each group owns one file, for
  example \texttt{players/group4/player.py}. It already contains a class
  (\texttt{Player4}) with one method to fill in:
\begin{lstlisting}
def act(self, snapshot: JellyfishSnapshot) -> Action:
    return Action(depth=snapshot.depth)
\end{lstlisting}

  \item \textbf{Understand when \texttt{act} runs.} Every day, the
  simulator calls \texttt{act} once for each living jellyfish in your
  group. Your job is to answer one question: \emph{which depth should this
  jellyfish ride today?}

  \item \textbf{Read what the jellyfish can sense.} The \texttt{snapshot}
  holds everything the jellyfish is allowed to know (see
  Section~\ref{sec:senses}). It does \textbf{not} know its $(x, y)$
  position or where any zone is. Navigating without that is the puzzle.

  \item \textbf{Remember things using memory.}
  \texttt{snapshot.memory} is a small array of bytes (8 by default, set
  with \texttt{-M}). Each byte holds a value from 0 to 255. Change it
  directly:
\begin{lstlisting}
snapshot.memory[0] = 5          # write
count = snapshot.memory[1]      # read
\end{lstlisting}
  This is the jellyfish's \textbf{only} memory. Do not store anything on
  \texttt{self} or in global variables between calls; the simulator does
  not even tell you which jellyfish is calling. You may not change the
  length of the memory. Offspring receive a copy of their parent's memory.

  \item \textbf{Return an action.}
\begin{lstlisting}
return Action(Depth.DEEP)
\end{lstlisting}
  Eating and reproducing happen automatically whenever the jellyfish is in
  the right zone in the right season, so depth is the only real decision.
  You can switch either off with \texttt{Action(depth, eat=False)} or
  \texttt{Action(depth, reproduce=False)}.

  \item \textbf{Use the provided randomness.} If you want random choices,
  use \texttt{self.rng} (a seeded \texttt{random.Random}) rather than
  Python's \texttt{random} module. Then the same seed always reproduces
  the same run.

  \item \textbf{Run and watch.}
\begin{lstlisting}[language=bash]
uv run main.py --players 4 --gui
uv run main.py --players 4 r e --gui   # compare with the baselines
\end{lstlisting}
  In VS Code you can also press F5 and pick a launch configuration.
\end{enumerate}

\subsection{A complete small example}

\begin{lstlisting}
from core.action import Action
from core.depth import Depth
from core.player import Player
from core.seasons import Season
from core.snapshot import JellyfishSnapshot


class Player4(Player):
    def act(self, s: JellyfishSnapshot) -> Action:
        if s.in_predator_zone:
            return Action(Depth.SURFACE)   # escape fast
        if s.in_feeding_zone and s.season == Season.SUMMER:
            return Action(Depth.DEEP)      # stay put and eat
        return Action(Depth.SHALLOW)       # otherwise drift
\end{lstlisting}

\subsection{Other things your class can use}

\begin{tabular}{@{}ll@{}}
\toprule
\texttt{self.memory\_size}   & number of bytes in \texttt{snapshot.memory} \\
\texttt{self.lifetime\_days} & $S \times 360$, the fixed lifespan of every jellyfish \\
\texttt{self.rng}            & seeded random number generator \\
\bottomrule
\end{tabular}

\medskip
You may add more files inside your own group folder.

% =====================================================================
\section{What a jellyfish senses}
\label{sec:senses}

\begin{tabular}{@{}p{0.42\textwidth}p{0.52\textwidth}@{}}
\toprule
\textbf{Field} & \textbf{Meaning} \\
\midrule
\texttt{day\_of\_year}, \texttt{year}, \texttt{season} & calendar (days 1--360; seasons spring, summer, fall, winter) \\
\texttt{birth\_day\_of\_year}, \texttt{birth\_year} & when this jellyfish was born \\
\texttt{age\_days} & convenience: days since birth \\
\texttt{depth} & its current depth \\
\texttt{heading}, \texttt{speed}, \texttt{velocity} & the current \emph{at its current depth} (see below) \\
\texttt{in\_feeding\_zone} & is it in the feeding zone? \\
\texttt{in\_breeding\_zone} & is it in the breeding zone? \\
\texttt{in\_predator\_zone} & is it in the predator zone? \\
\texttt{reproduction\_capacity} & offspring still allowed this winter \\
\texttt{memory} & its memory (a \texttt{bytearray}); mutate in place \\
\bottomrule
\end{tabular}

\subsection{Does it know the currents at the other depths?}

\textbf{Speeds: yes. Directions: no.}

\begin{itemize}[leftmargin=*]
  \item The \textbf{speeds} are fixed rules of the game, the same in every
  run: surface 40, shallow 20, medium 10 and deep 5~km/day. Your code can
  look them up at any time, e.g.\ \texttt{Depth.SURFACE.speed}.
  \item The \textbf{direction} of each current is chosen at random for
  every run. The snapshot only reports the current at the jellyfish's
  \emph{own} depth: its compass \texttt{heading} (0 = north, 90 = east),
  its \texttt{speed}, and its \texttt{velocity} as $(dx, dy)$ in km/day
  ($x$ grows east, $y$ grows south).
  \item To learn the direction at another depth, a jellyfish has to go
  there: next day the snapshot reports that depth's current. If you want
  to remember what you learned, you must store it in memory (for example,
  one byte per depth holding the heading divided by 360 and scaled to
  0--255). Since offspring inherit memory, knowledge can pass down the
  generations.
\end{itemize}

% =====================================================================
\section{The built-in players}

Three kinds of player ship with the simulator so you have something to
compare against. Select them with \texttt{-{}-players}.

\subsection{Random player (\texttt{r})}
Every day it picks one of the four depths uniformly at random (using
\texttt{self.rng}). It ignores the snapshot and never uses memory.
Its jellyfish wander aimlessly; this is the baseline to beat.

\subsection{Example player (\texttt{e})}
A short, readable strategy that uses only \textbf{one byte} of memory.
It checks these rules in order:
\begin{enumerate}[leftmargin=*]
  \item In the predator zone: ride the surface current (fastest) to get
  out quickly.
  \item In the feeding zone during summer: go deep (slowest current) to
  stay inside the zone and keep eating.
  \item In the breeding zone during winter with capacity left: go deep to
  stay and keep spawning.
  \item Otherwise \textbf{explore}: ride one current for 30 days, then
  pick a new random depth.
\end{enumerate}
To remember its exploration plan it packs two numbers into
\texttt{memory[0]}:
\begin{center}
\begin{tabular}{@{}lll@{}}
\toprule
bits 7--6 & exploration depth & 0--3 \\
bits 5--0 & days left on this leg & 0--63 \\
\bottomrule
\end{tabular}
\end{center}
It is not meant to be good; it shows how to read the snapshot and how to
pack information tightly into memory.

\subsection{Group stubs (\texttt{1}--\texttt{10})}
Until you change it, your player returns \texttt{Action(snapshot.depth)}:
it stays at its current depth forever. Jellyfish start at the surface, so
a stub drifts in a straight line at 40~km/day, wrapping around the ocean.

% =====================================================================
\section{Core features of the simulator}

\subsection{The world}
\begin{itemize}[leftmargin=*]
  \item A $1000 \times 1000$~km ocean that wraps around at every edge
  (a torus): leaving the right edge brings you back on the left, and
  leaving the top brings you back at the bottom.
  \item Three circular zones of radius 200~km: \textbf{feeding},
  \textbf{breeding} and \textbf{predator}. Feeding and breeding never
  overlap; the predator zone may overlap either.
  \item Four currents, one per depth, each with a fixed speed and a
  random direction per run.
  \item A 360-day year with four 90-day seasons: spring (days 1--90),
  summer (91--180), fall (181--270) and winter (271--360).
\end{itemize}

\subsection{What happens each day}
For every living jellyfish, in this order:
\begin{enumerate}[leftmargin=*]
  \item \textbf{Old age.} It dies once it has lived $S \times 360$ days.
  \item \textbf{Sensing.} Zone membership is checked at its position at
  the start of the day.
  \item \textbf{Predation.} If it starts the day in the predator zone, it
  is eaten with probability 5\%.
  \item \textbf{Decision.} The simulator calls your \texttt{act()}.
  \item \textbf{Eating.} In the feeding zone during summer, it eats. Its
  first 30 eating days (possibly spread over several summers) only grow
  it. After that, each eating day counts toward offspring.
  \item \textbf{Reproducing.} In the breeding zone during winter, with
  capacity left, it produces \emph{one} offspring at its location. The
  offspring copies the parent's memory and acts from the next day.
  \item \textbf{Drifting.} It moves to the chosen depth and drifts with
  that current for the day.
\end{enumerate}

On the first day of fall (day 181), each jellyfish's capacity is set to
\[
  \mbox{capacity} = \left\lfloor \frac{\mbox{eating days counted last summer}}{10} \right\rfloor
\]
and the counter resets. So reproduction always depends on the \emph{previous}
summer only.

\subsection{Start of a run}
$J$ jellyfish per group (set with \texttt{-J}) are born on day 1 of
year 1 at random spots in the breeding zone, at the surface, with memory
set to all zeros. Every group gets the same starting spots.

\subsection{Many groups in one ocean}
\texttt{-{}-players 1 2 3 \dots{} 10} runs every group in the same ocean at
the same time. Groups never interact, and each group has its own
predation dice, so one group's luck cannot affect another's. If one
group's code raises an error, returns something that is not an
\texttt{Action}, or changes the memory length, only that group stops; the
run continues for everyone else.

\subsection{Settings}

\begin{tabular}{@{}lll@{}}
\toprule
\textbf{Flag} & \textbf{Default} & \textbf{Meaning} \\
\midrule
\texttt{-{}-players} & \texttt{r} & \texttt{r} random, \texttt{e} example, \texttt{1}--\texttt{10} groups \\
\texttt{-{}-gui} & off & open the visualizer \\
\texttt{-{}-seed} & from clock & zones, current directions, predation, \texttt{self.rng} \\
\texttt{-T} & 3600 & days to simulate (360 days = 1 year) \\
\texttt{-M} & 8 & memory size in bytes (1--1024) \\
\texttt{-S} & 3 & lifetime in years (may be fractional) \\
\texttt{-J} & 10 & jellyfish each group starts with \\
\texttt{-{}-max\_population} & 2000 & per-group cap; extra births are skipped \\
\texttt{-{}-map\_path} & none & JSON file in \texttt{maps/} fixing zones and/or currents \\
\texttt{-{}-json} & off & print results as JSON only \\
\bottomrule
\end{tabular}

\subsection{Two ways to run}
\paragraph{Headless.} Prints a results table per group: final and peak
population, births, deaths from predators and from old age, the day the
group first reached $10J$ and $100J$, and how much time its code used.

\paragraph{GUI.} Shows the ocean with all three zones (feeding = solid
outline, breeding = dashed outline, predators = hatched) and the
jellyfish moving on it, coloured by group and drawn smaller the deeper
they are. The side panel has a season bar, a compass of the four
currents, a live population table and a population-over-time chart.
Eaten jellyfish play a short death animation.

\begin{center}
\begin{tabular}{@{}ll@{}}
\toprule
\textbf{Key} & \textbf{Action} \\
\midrule
\texttt{SPACE} & play / pause \\
$\rightarrow$ & advance one day \\
$\uparrow$ / $\downarrow$ & faster / slower \\
\texttt{L} & show one depth at a time \\
\texttt{T} & toggle trails \\
\texttt{G} & linear / log population chart \\
\texttt{1}--\texttt{9}, \texttt{0} & show / hide groups 1--10 \\
click & inspect a jellyfish (memory, age, eating counters) \\
\texttt{H} / \texttt{ESC} & help / quit \\
\bottomrule
\end{tabular}
\end{center}

\subsection{Code quality}
Pull requests must pass Ruff and the tests:
\begin{lstlisting}[language=bash]
uv run ruff format
uv run ruff check --fix
uv run python -m pytest
\end{lstlisting}

\end{document}
